\documentclass[10pt,a4paper]{amsart}
\usepackage[margin=19mm]{geometry}
\usepackage{amsmath,amssymb,mathtools}
\usepackage[scr=rsfs,cal=euler]{mathalfa}
\usepackage{xfrac}
\usepackage[unicode,hidelinks,bookmarks=false]{hyperref}
\newcommand{\ie}{i.e.\ }
\newcommand{\cf}{cf.\ }
\newcommand{\resp}{respectively\ }
\newcommand{\Subsection}{\subsection}
\providecommand{\nobreakdash}{\nobreak\hbox{-}}
\setlength{\emergencystretch}{2em}
\multlinegap0pt
\usepackage{amssymb}
\usepackage{mathtools}
\usepackage{bm}
\usepackage{xcolor}
\usepackage{algorithm}\usepackage{algpseudocode}
\newtheorem{proposition}{Proposition}
\title{Improving Cramér--Rao Bound And Its Variants: An Extrinsic Geometry Perspective}
\author{Sunder Ram Krishnan}
\date{Source: arXiv:2509.18978v2 (CC BY 4.0)}
\begin{document}
\maketitle

\noindent\emph{Source and licence.} Improving Cramér--Rao Bound And Its Variants: An Extrinsic Geometry Perspective. Version arXiv:2509.18978v2 (CC BY 4.0), distributed by arXiv under the Creative Commons Attribution 4.0 International licence, http://creativecommons.org/licenses/by/4.0/, as recorded in the arXiv licence metadata for this exact version and shown on the abstract page https://arxiv.org/abs/2509.18978v2. Full licence notice: \texttt{licenses/arxiv-2509.18978-CC-BY-4.0.txt}.\par\smallskip

\noindent\textit{Editorial scope.} Counted unit: the article's proposition prop:proj\_ineq (source0.tex lines 279--288). The article's complete original proof of it is delivered with it (source0.tex lines 290--315, one \texttt{proof} environment of 26 lines). No mathematics is written, repaired or re-derived: the statement and the proof are the article's own lines, in the article's own notation, and no retained line is substituted or re-flowed.  No further source-local context is needed: every cross-reference of every retained line resolves inside the two delivered spans.  Nothing was omitted from any retained span, so the package carries no omission record.  No retained line contains a citation, so the wrapper carries no bibliography and no bibliography entry.  No retained line contains a figure environment, an image-inclusion command or drawing code, so no image is delivered and the package declares no asset dependency.  Editorial formatting only, in addition: an AMS \texttt{amsart} preamble that restates the article's own declarations and macros used by the retained lines, namely the environments \texttt{proposition} on separate counters, together with the standard packages those lines need.  The retained environments print in this document as Proposition 1 in order of appearance, whereas the article prints the same items with the numbering of its own full document.  The complete native source of the article as distributed in the arXiv e-print for this exact version is bundled unchanged as \texttt{sources/185566/source0.tex}.
\begin{proposition}[Projection inequality]\label{prop:proj_ineq}
Let \(B\in\mathbb{R}^{d\times d}\) be the matrix with entries defined by the inner products of the projected lifted errors:
\[
B_{pq} := \big\langle \operatorname{Proj}_{\mathcal{T}_1}\widetilde Z^{(p)},\, \operatorname{Proj}_{\mathcal{T}_1}\widetilde Z^{(q)}\big\rangle,
\]
where \(\operatorname{Proj}_{\mathcal{T}_1}:H\to \mathcal{T}_1\) is the orthogonal projection. Then the estimator covariance satisfies:
\[
\Sigma(\theta) \succeq B.
\]
\end{proposition}

\begin{proof}
Fix an arbitrary direction vector \(v=(v_1,\dots,v_d)^\top\in\mathbb{R}^d\). Define the combined lifted error
\[
Z_v := \sum_{p=1}^d v_p\,\widetilde Z^{(p)} \in H.
\]
The squared norm of this vector in \(H\) corresponds to the variance in direction \(v\):
\[
\| Z_v \|^2 = \sum_{p,q} v_p v_q \langle \widetilde Z^{(p)}, \widetilde Z^{(q)} \rangle = v^\top\Sigma(\theta)v.
\]

Since the orthogonal projection onto a closed subspace of a Hilbert space is a contraction, we have:
\[
\|Z_v\|^2 \ge \big\| \operatorname{Proj}_{\mathcal{T}_1} Z_v\big\|^2.
\]
By the linearity of the projection operator,
\[
\operatorname{Proj}_{\mathcal{T}_1} Z_v = \sum_{p=1}^d v_p\,\operatorname{Proj}_{\mathcal{T}_1}\widetilde Z^{(p)}.
\]
The squared norm of the projected vector is then:
\[
\big\| \operatorname{Proj}_{\mathcal{T}_1} Z_v\big\|^2
= \sum_{p,q} v_p v_q\,\big\langle \operatorname{Proj}_{\mathcal{T}_1}\widetilde Z^{(p)},\operatorname{Proj}_{\mathcal{T}_1}\widetilde Z^{(q)}\big\rangle
= v^\top B v.
\]
Combining these results, we obtain \(v^\top\Sigma v \ge v^\top B v\) for all \(v\in\mathbb{R}^d\), which implies the matrix inequality \(\Sigma\succeq B\).
\end{proof}

\end{document}
